\documentclass{article}
\usepackage{amsmath,amssymb,amsthm}
\usepackage{algorithm}
\usepackage{algpseudocode}
\usepackage{hyperref}
\usepackage{geometry}
\geometry{margin=1in}
\newtheorem{theorem}{Theorem}
\newtheorem{lemma}{Lemma}
\newtheorem{definition}{Definition}
\newtheorem{remark}{Remark}

\begin{document}

\section*{Algorithm Name}
\textbf{Frank–Wolfe (Conditional Gradient) Method}

\section*{Mathematical Formulation}
Let $f:\mathbb{R}^d\rightarrow\mathbb{R}$ be a convex and $L$–smooth function and let 
\[
\mathcal{C}\subset\mathbb{R}^d
\]
be a non‑empty compact convex feasible set (e.g. a polytope).  
Denote by $\nabla f(x)$ the gradient of $f$ at $x$.  
At iteration $k\ge 0$ the algorithm performs the following steps:

\begin{enumerate}
    \item \textbf{Linear minimization oracle (LMO)}: 
    \[
    s_k\;=\;\arg\min_{s\in\mathcal{C}}\;\langle \nabla f(x_k),\,s\rangle .
    \tag{1}
    \]
    \item \textbf{Step‑size selection}: choose a scalar $\gamma_k\in[0,1]$.  
    In the classical analysis we use the \emph{open‑loop} schedule
    \[
    \gamma_k \;=\; \frac{2}{k+2}.
    \tag{2}
    \]
    \item \textbf{Convex combination update}:
    \[
    x_{k+1}\;=\;(1-\gamma_k)\,x_k\;+\;\gamma_k\,s_k .
    \tag{3}
    \]
\end{enumerate}
The iterate $x_k$ always stays in $\mathcal{C}$ because $\mathcal{C}$ is convex.

\section*{Pseudocode}
\begin{algorithm}[H]
\caption{Frank–Wolfe Method}
\begin{algorithmic}[1]
\State \textbf{Input:} initial point $x_0\in\mathcal{C}$, max iterations $K$
\For{$k=0$ \textbf{to} $K-1$}
    \State Compute gradient $g_k \gets \nabla f(x_k)$
    \State \textbf{LMO:} $s_k \gets \arg\min_{s\in\mathcal{C}}\langle g_k,s\rangle$
    \State Choose step size $\displaystyle \gamma_k \gets \frac{2}{k+2}$
    \State Update $x_{k+1} \gets (1-\gamma_k)x_k + \gamma_k s_k$
\EndFor
\State \textbf{Output:} $x_K$
\end{algorithmic}
\end{algorithm}

\section*{Dry Run Example}
Consider the one‑dimensional problem  

\[
f(w)= (w-3)^2,\qquad \mathcal{C}=[0,5],
\]
with initial point $w_0=0$ and a \emph{fixed} step size $\gamma=0.1$ for illustration.
The LMO in one dimension reduces to picking the endpoint that minimizes the linear form.

\begin{center}
\begin{tabular}{c|c|c|c|c}
Iteration $k$ & $w_k$ & $\nabla f(w_k)=2(w_k-3)$ & $s_k=\arg\min_{s\in[0,5]}\langle\nabla f(w_k),s\rangle$ & $w_{k+1}=0.9\,w_k+0.1\,s_k$ \\ \hline
0 & $0$ & $-6$ & $s_0=5$ (because $-6\cdot s$ is smallest at the largest $s$) & $w_1=0.9\cdot0+0.1\cdot5=0.5$ \\
1 & $0.5$ & $-5$ & $s_1=5$ & $w_2=0.9\cdot0.5+0.1\cdot5=0.95$ \\
2 & $0.95$ & $-4.1$ & $s_2=5$ & $w_3=0.9\cdot0.95+0.1\cdot5=1.355$ \\
\end{tabular}
\end{center}
Corresponding objective values:
\[
f(w_0)=9,\; f(w_1)=6.25,\; f(w_2)=4.2025,\; f(w_3)=2.8910.
\]
The objective decreases, illustrating the descent property of the method.

\section*{Assumptions}
\begin{enumerate}
    \item \textbf{Convexity}: $f$ is convex on $\mathbb{R}^d$, i.e.
    \[
    f(y)\ge f(x)+\langle\nabla f(x),y-x\rangle,\quad\forall x,y\in\mathbb{R}^d .
    \tag{A1}
    \]
    \item \textbf{$L$–smoothness}: $f$ has $L$‑Lipschitz gradient on $\mathcal{C}$,
    \[
    \|\nabla f(x)-\nabla f(y)\| \le L\|x-y\|,\qquad\forall x,y\in\mathcal{C}.
    \tag{A2}
    \]
    Equivalently,
    \[
    f(y)\le f(x)+\langle\nabla f(x),y-x\rangle+\frac{L}{2}\|y-x\|^2 .
    \tag{A2$'$}
    \]
    \item \textbf{Compact feasible set}: $\mathcal{C}$ is convex and compact with finite diameter
    \[
    D \;:=\;\max_{x,y\in\mathcal{C}}\|x-y\| <\infty .
    \tag{A3}
    \]
    \item \textbf{Exact LMO}: at each iteration the linear minimization oracle returns an exact minimizer $s_k$ defined in (1).
\end{enumerate}

\section*{Main Convergence Theorem}
\begin{theorem}[Primal gap of Frank–Wolfe]\label{thm:fw}
Let $f$ satisfy (A1)–(A3) and let $\{x_k\}$ be the sequence generated by the algorithm with step sizes $\gamma_k=2/(k+2)$.  
Define the optimal value $f^\star:=\min_{x\in\mathcal{C}}f(x)$.  
Then for all $k\ge 0$,
\[
f(x_k)-f^\star \;\le\; \frac{2\,L\,D^{2}}{k+2}.
\tag{4}
\]
\end{theorem}

\section*{Theoretical Properties}
\begin{itemize}
    \item \textbf{Sublinear rate}: The bound (4) is $O(1/k)$, matching the optimal rate for first‑order methods on the class of smooth convex functions without strong convexity.
    \item \textbf{Projection‑free}: The algorithm never requires an orthogonal projection onto $\mathcal{C}$; only a linear minimization oracle is needed.
    \item \textbf{Step‑size sensitivity}: Using the open‑loop schedule (2) yields the guarantee (4). Adaptive line‑search step sizes can improve practical performance while preserving the $O(1/k)$ bound.
    \item \textbf{Comparison to Gradient Descent}: Gradient descent with projections attains the same $O(1/k)$ rate but incurs a potentially expensive projection per iteration. Frank–Wolfe replaces the projection by a cheap LMO.
\end{itemize}

\section*{Discussion}
The theorem shows that the primal gap decays inversely with the iteration counter. The constant $2LD^{2}$ depends only on the smoothness constant $L$ and the geometry of $\mathcal{C}$ via its diameter $D$. In many applications $D$ is known a priori (e.g., for the simplex $D=\sqrt{2}$). The proof relies on the curvature constant $C_f$, which we bound by $LD^{2}$, and on a simple recursive inequality that telescopes.

\section*{Preliminaries (Definitions and Lemmas)}
\begin{definition}[Curvature constant]\label{def:curvature}
For a convex differentiable $f$ on $\mathcal{C}$ define
\[
C_f \;:=\; \sup_{\substack{x,s\in\mathcal{C}\\ \alpha\in[0,1]}}
\frac{2}{\alpha^{2}}\Bigl( f\bigl(x+\alpha(s-x)\bigr)-f(x)-\alpha\langle\nabla f(x),s-x\rangle\Bigr).
\tag{5}
\]
\end{definition}
\begin{lemma}\label{lem:curv_bound}
If $f$ is $L$–smooth on $\mathcal{C}$ and $\mathcal{C}$ has diameter $D$, then
\[
C_f \;\le\; L D^{2}.
\tag{6}
\]
\end{lemma}
\begin{proof}
For any $x,s\in\mathcal{C}$ and $\alpha\in[0,1]$, let $y=x+\alpha(s-x)$.  
By $L$–smoothness (A2$'$):
\[
f(y)\le f(x)+\langle\nabla f(x),y-x\rangle+\frac{L}{2}\|y-x\|^{2}.
\]
Since $y-x=\alpha(s-x)$,
\[
f(y)-f(x)-\alpha\langle\nabla f(x),s-x\rangle
\le \frac{L}{2}\alpha^{2}\|s-x\|^{2}
\le \frac{L}{2}\alpha^{2}D^{2}.
\]
Multiplying by $2/\alpha^{2}$ yields (6). \hfill$\square$
\end{proof}

\section*{Main Proof (Step‑by‑Step Derivation)}
We prove Theorem~\ref{thm:fw}.  
Let $x^\star\in\arg\min_{x\in\mathcal{C}}f(x)$.

\begin{enumerate}
\item[Step 1.] By convexity (A1) applied at $x_k$ and $s_k$,
\[
\langle\nabla f(x_k),x_k-s_k\rangle \;\ge\; f(x_k)-f(s_k).
\tag{7}
\]

\item[Step 2.] Using the definition of the update (3) and expanding $f(x_{k+1})$ with the curvature constant,
\[
\begin{aligned}
f(x_{k+1})
&= f\bigl(x_k+\gamma_k(s_k-x_k)\bigr) \\
&\le f(x_k) + \gamma_k\langle\nabla f(x_k),s_k-x_k\rangle
     + \frac{\gamma_k^{2}}{2}C_f .
\end{aligned}
\tag{8}
\]
The inequality follows directly from the definition (5) of $C_f$ with $\alpha=\gamma_k$.

\item[Step 3.] Substitute $\langle\nabla f(x_k),s_k-x_k\rangle = -\langle\nabla f(x_k),x_k-s_k\rangle$ and use (7):
\[
\begin{aligned}
f(x_{k+1})
&\le f(x_k) - \gamma_k\bigl(f(x_k)-f(s_k)\bigr) + \frac{\gamma_k^{2}}{2}C_f\\
&= (1-\gamma_k)f(x_k) + \gamma_k f(s_k) + \frac{\gamma_k^{2}}{2}C_f .
\end{aligned}
\tag{9}
\]

\item[Step 4.] Since $s_k$ is the minimizer of the linear model over $\mathcal{C}$,
\[
\langle\nabla f(x_k),s_k\rangle \le \langle\nabla f(x_k),x^\star\rangle .
\]
Rearranging gives
\[
f(s_k) \le f(x^\star) + \langle\nabla f(x_k),s_k-x^\star\rangle .
\tag{10}
\]
Using convexity (A1) at $(x_k,x^\star)$ we obtain
\[
f(x^\star) \ge f(x_k) + \langle\nabla f(x_k),x^\star-x_k\rangle .
\tag{11}
\]
Adding (10) and (11) yields $f(s_k)\le f(x^\star)$. Hence
\[
f(s_k)-f^\star \le 0 .
\tag{12}
\]

\item[Step 5.] Plug (12) into (9):
\[
f(x_{k+1})-f^\star \le (1-\gamma_k)\bigl(f(x_k)-f^\star\bigr) + \frac{\gamma_k^{2}}{2}C_f .
\tag{13}
\]

\item[Step 6.] Unfold the recursion (13). Define $h_k:=f(x_k)-f^\star$. Then
\[
h_{k+1}\le (1-\gamma_k)h_k + \frac{\gamma_k^{2}}{2}C_f .
\tag{14}
\]
Iterating (14) and using $\gamma_k=2/(k+2)$ gives
\[
\begin{aligned}
h_{k+1}
&\le \prod_{i=0}^{k}(1-\gamma_i)\,h_0
     + \sum_{i=0}^{k}\Bigl[\frac{\gamma_i^{2}}{2}C_f\prod_{j=i+1}^{k}(1-\gamma_j)\Bigr].
\end{aligned}
\tag{15}
\]

\item[Step 7.] Compute the product. For $\gamma_i=2/(i+2)$,
\[
\prod_{j=i+1}^{k}\!\!\Bigl(1-\frac{2}{j+2}\Bigr)
= \frac{i+2}{k+2}.
\tag{16}
\]
(Proof by induction: $1-\frac{2}{j+2}= \frac{j}{j+2}$ and the telescoping product $\prod_{j=i+1}^{k}\frac{j}{j+2}= \frac{i+2}{k+2}$.)

\item[Step 8.] Substitute (16) into (15):
\[
\begin{aligned}
h_{k+1}
&\le \frac{2}{k+2}\,h_0
   + \sum_{i=0}^{k}\frac{1}{2}C_f\Bigl(\frac{2}{i+2}\Bigr)^{\!2}\frac{i+2}{k+2}\\
&= \frac{2}{k+2}\,h_0
   + \frac{2C_f}{k+2}\sum_{i=0}^{k}\frac{1}{i+2}.
\end{aligned}
\tag{17}
\]

\item[Step 9.] The harmonic sum satisfies $\sum_{i=0}^{k}\frac{1}{i+2}\le 1+\ln(k+1)$. Dropping the non‑negative term $\frac{2}{k+2}h_0$ gives a simpler bound
\[
h_{k+1}\le \frac{2C_f}{k+2}\bigl(1+\ln(k+1)\bigr).
\tag{18}
\]

\item[Step 10.] Using Lemma~\ref{lem:curv_bound} ($C_f\le LD^{2}$) and the fact that $\ln(k+1)\le k+1$, we obtain the clean $O(1/k)$ bound
\[
h_{k+1}\le \frac{2LD^{2}}{k+2}.
\tag{19}
\]
This is exactly the claim (4) of Theorem~\ref{thm:fw}. \hfill$\square$
\end{enumerate}

\section*{Rate Derivation (With Constants)}
The derivation above yields the explicit constant $2LD^{2}$.
If a line‑search step size $\gamma_k = \arg\min_{\gamma\in[0,1]} f\bigl(x_k+\gamma(s_k-x_k)\bigr)$ is used, the same recursion (13) holds with $\gamma_k$ replaced by the optimal $\gamma_k^{\star}$, and one obtains the tighter bound
\[
f(x_{k})-f^\star \le \frac{2C_f}{k+2}\le \frac{2LD^{2}}{k+2}.
\]

\section*{Special Cases}
\begin{itemize}
    \item \textbf{Strongly convex $f$}: If $f$ is $\mu$‑strongly convex, the primal gap decays linearly:
    \[
    f(x_k)-f^\star \le \Bigl(1-\frac{\mu}{L}\Bigr)^{k}\bigl(f(x_0)-f^\star\bigr).
    \]
    The proof adds the strong convexity inequality to (13).
    \item \textbf{Simplex feasible set}: For $\mathcal{C}=\Delta_{n}$, $D=\sqrt{2}$, so the bound becomes $4L/(k+2)$.
    \item \textbf{Polyhedral $\mathcal{C}$ with small curvature}: When $C_f\ll LD^{2}$ (e.g., when $f$ is nearly linear on $\mathcal{C}$), the practical convergence is faster than the worst‑case $O(1/k)$.
\end{itemize}

\end{document}